\section{Research Background}

The rapid growth of the RegTech industry is driven by the increasing complexity, volume, and frequency of regulatory requirements. Regulatory information is often contained in large collections of semi-structured and semi-labeled documents, where relevant knowledge is distributed across text, metadata, references, and hierarchical structures. Traditional compliance systems face challenges in efficiently retrieving and connecting this heterogeneous information.

A fundamental challenge is the interconnected nature of regulatory knowledge. Regulations frequently reference other regulations, standards, technical specifications, and exceptions, forming complex networks of dependencies. As a result, understanding a specific requirement often requires identifying relationships between entities and retrieving information across multiple sources. Conventional keyword-based and document-level retrieval methods are insufficient for capturing these semantic connections.

Recent advances in artificial intelligence, including Transformer-based models, Retrieval-Augmented Generation (RAG), and associative memory architectures, provide new opportunities for representing and retrieving interconnected regulatory knowledge. This thesis investigates approaches for improving regulatory information retrieval through semantic representation and associative retrieval methods, with a focus on entity-aware knowledge retrieval in RegTech applications.
